% Shared preamble for the course decision diagrams.
% Each diagram file starts with % Shared preamble for the course decision diagrams.
% Each diagram file starts with % Shared preamble for the course decision diagrams.
% Each diagram file starts with \input{diagram_style} and contains one tikzpicture.
% Compile with pdflatex; build_diagrams.py then crops each page to its content.
\documentclass[11pt]{article}
\usepackage[a3paper,landscape,margin=10mm]{geometry}
\usepackage[T1]{fontenc}
\usepackage{tgheros}                 % Neudata body font (as in neudata.cls)
\renewcommand{\sfdefault}{qhv}
\renewcommand{\familydefault}{\sfdefault}
\usepackage[scaled=0.95,varqu]{zi4}
\usepackage{tikz}
\usetikzlibrary{arrows.meta,positioning,fit,calc}
\pagestyle{empty}
\setlength{\parindent}{0pt}
\hyphenpenalty=10000
\exhyphenpenalty=10000

% Neudata palette (neu-data/.github BRAND.md); names kept so the diagram sources need no edits
\definecolor{cblue}{HTML}{0B376C}    % Neudata Blue: question boxes, links between diagrams
\definecolor{corange}{HTML}{B5652B}  % muted contrasting colour: warnings
\definecolor{cgreen}{HTML}{055F56}   % Neudata Teal: method boxes
\definecolor{clight}{HTML}{EAF2F5}   % Neudata Mist: question box fill
\definecolor{ctext}{HTML}{1F2D33}    % Neudata Ink: text
\definecolor{cgrey}{HTML}{8A99A3}    % Neudata Grey: further-learning routes
\definecolor{cpale}{HTML}{F4F8FA}
\definecolor{cnavy}{HTML}{04242F}    % Deep Navy: titles

\tikzset{
  every node/.style={text=ctext, font=\small},
  title/.style={font=\Large\bfseries, text=cnavy, anchor=west},
  subtitle/.style={font=\normalsize, text=ctext!72, anchor=west},
  q/.style={draw=cblue, fill=clight, rounded corners=3pt, line width=0.8pt,
            align=center, inner sep=6pt, text width=#1},
  q/.default=5cm,
  m/.style={draw=cgreen, fill=white, rounded corners=3pt, line width=1pt,
            align=left, inner sep=6pt, text width=#1},
  m/.default=5cm,
  opt/.style={draw=cblue, fill=white, rounded corners=3pt, line width=0.6pt,
              align=center, inner sep=5pt, text width=#1},
  opt/.default=3.2cm,
  further/.style={draw=cgrey, dashed, fill=white, rounded corners=3pt, line width=0.8pt,
                  align=left, inner sep=6pt, text=ctext, text width=#1},
  further/.default=5cm,
  warn/.style={draw=corange, fill=white, rounded corners=3pt, line width=1pt,
               align=left, inner sep=6pt, text width=#1},
  warn/.default=6cm,
  note/.style={align=left, text width=#1, font=\footnotesize, text=ctext},
  note/.default=5cm,
  band/.style={font=\normalsize\bfseries, text=cgreen, anchor=west, align=left},
  flow/.style={-{Stealth[length=2.6mm]}, line width=0.8pt, draw=ctext!75},
  flowg/.style={-{Stealth[length=2.6mm]}, line width=0.8pt, draw=cgrey, dashed},
  lab/.style={font=\footnotesize, fill=white, inner sep=2pt, text=ctext},
  code/.style={font=\ttfamily\footnotesize},
}
\newcommand{\rc}[1]{{\ttfamily\footnotesize #1}}
\newcommand{\dref}[1]{{\color{cblue}#1}}
 and contains one tikzpicture.
% Compile with pdflatex; build_diagrams.py then crops each page to its content.
\documentclass[11pt]{article}
\usepackage[a3paper,landscape,margin=10mm]{geometry}
\usepackage[T1]{fontenc}
\usepackage{tgheros}                 % Neudata body font (as in neudata.cls)
\renewcommand{\sfdefault}{qhv}
\renewcommand{\familydefault}{\sfdefault}
\usepackage[scaled=0.95,varqu]{zi4}
\usepackage{tikz}
\usetikzlibrary{arrows.meta,positioning,fit,calc}
\pagestyle{empty}
\setlength{\parindent}{0pt}
\hyphenpenalty=10000
\exhyphenpenalty=10000

% Neudata palette (neu-data/.github BRAND.md); names kept so the diagram sources need no edits
\definecolor{cblue}{HTML}{0B376C}    % Neudata Blue: question boxes, links between diagrams
\definecolor{corange}{HTML}{B5652B}  % muted contrasting colour: warnings
\definecolor{cgreen}{HTML}{055F56}   % Neudata Teal: method boxes
\definecolor{clight}{HTML}{EAF2F5}   % Neudata Mist: question box fill
\definecolor{ctext}{HTML}{1F2D33}    % Neudata Ink: text
\definecolor{cgrey}{HTML}{8A99A3}    % Neudata Grey: further-learning routes
\definecolor{cpale}{HTML}{F4F8FA}
\definecolor{cnavy}{HTML}{04242F}    % Deep Navy: titles

\tikzset{
  every node/.style={text=ctext, font=\small},
  title/.style={font=\Large\bfseries, text=cnavy, anchor=west},
  subtitle/.style={font=\normalsize, text=ctext!72, anchor=west},
  q/.style={draw=cblue, fill=clight, rounded corners=3pt, line width=0.8pt,
            align=center, inner sep=6pt, text width=#1},
  q/.default=5cm,
  m/.style={draw=cgreen, fill=white, rounded corners=3pt, line width=1pt,
            align=left, inner sep=6pt, text width=#1},
  m/.default=5cm,
  opt/.style={draw=cblue, fill=white, rounded corners=3pt, line width=0.6pt,
              align=center, inner sep=5pt, text width=#1},
  opt/.default=3.2cm,
  further/.style={draw=cgrey, dashed, fill=white, rounded corners=3pt, line width=0.8pt,
                  align=left, inner sep=6pt, text=ctext, text width=#1},
  further/.default=5cm,
  warn/.style={draw=corange, fill=white, rounded corners=3pt, line width=1pt,
               align=left, inner sep=6pt, text width=#1},
  warn/.default=6cm,
  note/.style={align=left, text width=#1, font=\footnotesize, text=ctext},
  note/.default=5cm,
  band/.style={font=\normalsize\bfseries, text=cgreen, anchor=west, align=left},
  flow/.style={-{Stealth[length=2.6mm]}, line width=0.8pt, draw=ctext!75},
  flowg/.style={-{Stealth[length=2.6mm]}, line width=0.8pt, draw=cgrey, dashed},
  lab/.style={font=\footnotesize, fill=white, inner sep=2pt, text=ctext},
  code/.style={font=\ttfamily\footnotesize},
}
\newcommand{\rc}[1]{{\ttfamily\footnotesize #1}}
\newcommand{\dref}[1]{{\color{cblue}#1}}
 and contains one tikzpicture.
% Compile with pdflatex; build_diagrams.py then crops each page to its content.
\documentclass[11pt]{article}
\usepackage[a3paper,landscape,margin=10mm]{geometry}
\usepackage[T1]{fontenc}
\usepackage{tgheros}                 % Neudata body font (as in neudata.cls)
\renewcommand{\sfdefault}{qhv}
\renewcommand{\familydefault}{\sfdefault}
\usepackage[scaled=0.95,varqu]{zi4}
\usepackage{tikz}
\usetikzlibrary{arrows.meta,positioning,fit,calc}
\pagestyle{empty}
\setlength{\parindent}{0pt}
\hyphenpenalty=10000
\exhyphenpenalty=10000

% Neudata palette (neu-data/.github BRAND.md); names kept so the diagram sources need no edits
\definecolor{cblue}{HTML}{0B376C}    % Neudata Blue: question boxes, links between diagrams
\definecolor{corange}{HTML}{B5652B}  % muted contrasting colour: warnings
\definecolor{cgreen}{HTML}{055F56}   % Neudata Teal: method boxes
\definecolor{clight}{HTML}{EAF2F5}   % Neudata Mist: question box fill
\definecolor{ctext}{HTML}{1F2D33}    % Neudata Ink: text
\definecolor{cgrey}{HTML}{8A99A3}    % Neudata Grey: further-learning routes
\definecolor{cpale}{HTML}{F4F8FA}
\definecolor{cnavy}{HTML}{04242F}    % Deep Navy: titles

\tikzset{
  every node/.style={text=ctext, font=\small},
  title/.style={font=\Large\bfseries, text=cnavy, anchor=west},
  subtitle/.style={font=\normalsize, text=ctext!72, anchor=west},
  q/.style={draw=cblue, fill=clight, rounded corners=3pt, line width=0.8pt,
            align=center, inner sep=6pt, text width=#1},
  q/.default=5cm,
  m/.style={draw=cgreen, fill=white, rounded corners=3pt, line width=1pt,
            align=left, inner sep=6pt, text width=#1},
  m/.default=5cm,
  opt/.style={draw=cblue, fill=white, rounded corners=3pt, line width=0.6pt,
              align=center, inner sep=5pt, text width=#1},
  opt/.default=3.2cm,
  further/.style={draw=cgrey, dashed, fill=white, rounded corners=3pt, line width=0.8pt,
                  align=left, inner sep=6pt, text=ctext, text width=#1},
  further/.default=5cm,
  warn/.style={draw=corange, fill=white, rounded corners=3pt, line width=1pt,
               align=left, inner sep=6pt, text width=#1},
  warn/.default=6cm,
  note/.style={align=left, text width=#1, font=\footnotesize, text=ctext},
  note/.default=5cm,
  band/.style={font=\normalsize\bfseries, text=cgreen, anchor=west, align=left},
  flow/.style={-{Stealth[length=2.6mm]}, line width=0.8pt, draw=ctext!75},
  flowg/.style={-{Stealth[length=2.6mm]}, line width=0.8pt, draw=cgrey, dashed},
  lab/.style={font=\footnotesize, fill=white, inner sep=2pt, text=ctext},
  code/.style={font=\ttfamily\footnotesize},
}
\newcommand{\rc}[1]{{\ttfamily\footnotesize #1}}
\newcommand{\dref}[1]{{\color{cblue}#1}}

\begin{document}
\begin{tikzpicture}[x=1cm,y=1cm]
\node[title] at (0,2.2) {D2. Comparing a continuous or ordinal outcome between groups};
\node[subtitle] at (0,1.5) {Report the size of the difference with a 95\% CI. The test is only one part of the answer.};

\node[q=11cm, anchor=north] (root) at (12.6,0.6) {\textbf{Continuous or ordinal outcome compared between groups}\\\footnotesize e.g.\ SBP, CRP or satisfaction by sex, smoking, district or time};
\node[q=7cm, below=6mm of root] (design) {\textbf{Are the observations independent or paired?}};
\draw[flow] (root) -- (design);

% ---- independent branch
\node[q=6cm, anchor=north] (ngroups) at ($(design.south)+(-4.3,-1.0)$) {\textbf{Independent groups: how many?}};
\draw[flow] (design.south) -- ++(0,-0.45) -| node[lab, pos=0.25] {independent} (ngroups.north);

\node[opt=4cm, anchor=north] (two) at ($(ngroups.south)+(-4.3,-0.9)$) {\textbf{Two groups}};
\node[opt=4cm, anchor=north] (three) at ($(ngroups.south)+(4.3,-0.9)$) {\textbf{Three or more groups}};
\draw[flow] (ngroups.south) -- ++(0,-0.4) -| (two.north);
\draw[flow] (ngroups.south) -- ++(0,-0.4) -| (three.north);

\node[m=7.6cm, below=6mm of two] (welch) {\textbf{Welch t-test} (sensible default)\\
Report: difference in means, 95\% CI\\
\rc{t.test(sbp\_0 \textasciitilde{} sex, data = hearts)}\\
Student's t-test (\rc{var.equal = TRUE}) only if equal spread is justified.};
\node[m=7.6cm, below=4mm of welch] (logt) {\textbf{Skewed, positive outcome: t-test on the log scale}\\
Report: ratio of geometric means, 95\% CI\\
\rc{t.test(log(crp) \textasciitilde{} smoker, data = hearts)}};
\node[m=7.6cm, below=4mm of logt] (wrs) {\textbf{Wilcoxon rank-sum} (Mann-Whitney)\\
For ordinal outcomes, or when ranks suit the question. Asks whether values in one group tend to be higher.\\
\rc{wilcox.test(crp \textasciitilde{} smoker, data = hearts, conf.int = TRUE)}\\
Estimate: Hodges-Lehmann shift, not a difference in medians.};

\node[m=7.6cm, below=6mm of three] (anova) {\textbf{Welch ANOVA} (default) \textbf{or classical ANOVA}\\
Report: group means with 95\% CIs\\
\rc{oneway.test(sbp\_0 \textasciitilde{} district, data = hearts)}\\
\rc{fit\_aov <- aov(sbp\_0 \textasciitilde{} district, data = hearts)}};
\node[m=7.6cm, below=4mm of anova] (kw) {\textbf{Kruskal-Wallis test}\\
Rank-based; ordinal outcomes or when ranks suit the question\\
\rc{kruskal.test(crp \textasciitilde{} district, data = hearts)}};
\node[warn=7.6cm, below=8mm of kw] (post) {\textbf{An overall test does not say which groups differ.}\\
If needed, compare pairs and adjust for multiple comparisons:\\
\rc{pairwise.t.test(..., pool.sd = FALSE,}\\
\rc{\ \ p.adjust.method = "holm")}\quad or \quad\rc{TukeyHSD(fit\_aov)}};
\draw[flow] (kw.south) -- (post.north);

% ---- paired branch
\node[q=6.6cm, anchor=north] (paired) at ($(design.south)+(8.7,-1.0)$) {\textbf{Paired: analyse the within-pair differences}\\\footnotesize same people measured twice, or matched pairs};
\draw[flow] (design.south) -- ++(0,-0.45) -| node[lab, pos=0.25] {paired} (paired.north);
\node[m=6.6cm, below=6mm of paired] (pt) {\textbf{Paired t-test}\\
Report: mean difference, 95\% CI\\
\rc{t.test(prog\$sbp\_12, prog\$sbp\_0,}\\\rc{\ \ \ \ \ \ \ paired = TRUE)}};
\node[m=6.6cm, below=4mm of pt] (sr) {\textbf{Wilcoxon signed-rank test}\\
Rank-based alternative for the differences\\
\rc{wilcox.test(..., paired = TRUE)}};
\node[warn=6.6cm, below=4mm of sr] (wrong) {\textbf{Do not use an independent-groups test on paired data.} It ignores that each person is compared with themselves (wider, wrong CI).};
\node[further=6.6cm, below=4mm of wrong] (rep) {Three or more time points per person: see \dref{D6} (repeated measures).};

% ---- assumptions box
\node[warn=25.3cm, anchor=north west] at ($(wrs.south -| 0.2,0)+(0,-0.7)$) {\textbf{Assumptions and choices.}
t-tests and ANOVA assume independent observations (from the design) and approximately normal sample means. In large groups this holds even with moderate skew; with small groups, strong skew or outliers matter more. Check residual plots.
Rank-based tests are not assumption-free: they also need independent observations, and they are not in general tests of medians.
Do not choose a test from sample size alone or from a normality-test p-value: look at the data and decide which summary answers your question.};
\end{tikzpicture}
\end{document}
